\documentclass[11pt,a4paper]{article}

\usepackage[spanish,es-nodecimaldot]{babel}
\usepackage[utf8]{inputenc}
\usepackage[T1]{fontenc}
\usepackage{lmodern}
\usepackage{geometry}
\usepackage{microtype}
\usepackage{amsmath,amssymb,amsfonts,mathtools}
\usepackage{enumitem}
\usepackage{xcolor}
\usepackage{hyperref}
\usepackage{fancyhdr}

\geometry{margin=2.35cm}
\setlength{\headheight}{14pt}
\setlength{\parindent}{0pt}
\setlength{\parskip}{0.55em}
\setlist[enumerate]{leftmargin=2.2em,itemsep=0.55em,topsep=0.45em}

\definecolor{repblue}{HTML}{12324A}
\definecolor{repteal}{HTML}{2A6F6B}
\definecolor{repgray}{HTML}{5C6770}
\definecolor{repline}{HTML}{C8D2DA}
\definecolor{replight}{HTML}{F4F7F9}

\hypersetup{
  colorlinks=true,
  linkcolor=repblue,
  urlcolor=repteal,
  pdftitle={TRGF Padilla Tarea 3},
  pdfauthor={Juan Ignacio Padilla Barrientos}
}

\pagestyle{fancy}
\fancyhf{}
\fancyhead[L]{\textcolor{repblue}{TRGF Padilla}}
\fancyhead[R]{\textcolor{repgray}{Tarea 3}}
\fancyfoot[C]{\thepage}
\renewcommand{\headrulewidth}{0.3pt}
\renewcommand{\footrulewidth}{0pt}

\newcommand{\CC}{\mathbb{C}}
\newcommand{\ZZ}{\mathbb{Z}}

\newcommand{\topbanner}{%
  {\color{repblue}\rule{\linewidth}{1.1pt}}\par
  \vspace{0.9em}
  \begin{center}
    {\LARGE\bfseries\color{repblue} TRGF Padilla -- Tarea 3}\par
    \vspace{0.35em}
    {\large\color{repteal} Seminario UCR I 2026}\par
    \vspace{0.55em}
  \end{center}
  \vspace{0.55em}
  {\color{repblue}\rule{\linewidth}{0.8pt}}\par
}

\newcommand{\problem}[2]{%
  \par\vspace{1.1em}
  {\large\bfseries\color{repblue} Problema #1}\par
  {\color{repteal}\rule{\linewidth}{0.45pt}}\par\smallskip
  #2
}

\newcommand{\worklines}[1]{%
  \par\vspace{0.35em}
  {\color{repline}\rule{\linewidth}{0.4pt}}\par
  \begingroup
  \count0=1
  \loop
    \ifnum\count0<#1
      \vspace{0.72cm}
      {\color{repline}\rule{\linewidth}{0.4pt}}\par
      \advance\count0 by 1
  \repeat
  \endgroup
  \vspace{0.35em}
}

\begin{document}

\color{black}
\topbanner

{\small\color{repgray}
}

\problem{1}{%
\textit{\textquestiondown Para cuales $n$ se cumple que hay un homomorfismo sobreyectivo de grupos $S_{n+1} \to S_n$?}
}

\textbf{Respuesta.} Existe un homomorfismo sobreyectivo $S_{n+1}\to S_n$ si y solo si $n\in\{1,2,3\}$.

\medskip
\textbf{Hecho que usaremos.}

\textit{Primer teorema de isomorfismo:} si $\phi\colon G\to H$ es un homomorfismo sobreyectivo de grupos, entonces $G/\ker\phi\cong H$. En particular, $|\ker\phi|=|G|/|H|$.

\medskip
\textbf{Lema 1 (Jordan).} Si $m\ge 5$, entonces $A_m$ es simple.

\medskip
\textbf{Caso $n=1$.}
$S_2\to S_1=\{e\}$. El unico homomorfismo (el trivial) es sobreyectivo.

\medskip
\textbf{Caso $n=2$.}
El homomorfismo signo $\operatorname{sgn}\colon S_3\to\{+1,-1\}$ es sobreyectivo.
Como $\{+1,-1\}\cong \ZZ/2\ZZ\cong S_2$, obtenemos el homomorfismo sobreyectivo deseado.

\medskip
\textbf{Caso $n=3$.}
Considere las tres particiones de $\{1,2,3,4\}$ en dos pares:
\[
  P_1=\bigl\{\{1,2\},\{3,4\}\bigr\},\quad
  P_2=\bigl\{\{1,3\},\{2,4\}\bigr\},\quad
  P_3=\bigl\{\{1,4\},\{2,3\}\bigr\}.
\]
Estas son \emph{todas} las particiones de $\{1,2,3,4\}$ en dos pares
(hay $\binom{4}{2}/2!=3$).
Cada $\sigma\in S_4$ actua sobre una particion aplicando $\sigma$ a cada elemento;
como $\sigma$ es biyeccion, el resultado es otra particion en dos pares,
y como solo hay tres, $\sigma$ permuta $\{P_1,P_2,P_3\}$; ademas, $(1\;2)$ induce la trasposicion $(P_2\;P_3)$ y $(1\;2\;3)$ induce el $3$-ciclo $(P_1\;P_3\;P_2)$.

\textit{Que $\phi$ es homomorfismo:} esto es inmediato, pues una accion de grupo sobre un conjunto $X$ es precisamente un homomorfismo $G\to S(X)$; aqui $X=\{P_1,P_2,P_3\}$ y $S(X)\cong S_3$.

\textit{Sobreyectividad:}
la imagen de $\phi$ contiene una trasposicion y un $3$-ciclo en $S_3$.
En efecto, $\phi((1\;2))=(P_2\;P_3)$ y $\phi((1\;2\;3))=(P_1\;P_3\;P_2)$.
Un subgrupo de $S_3$ que contiene una trasposicion y un $3$-ciclo es todo $S_3$
(pues $|S_3|=6$ y el subgrupo generado tiene orden divisible por $2$ y por $3$, asi que es $6$).
Por tanto $\phi$ es sobreyectivo.

\textit{Nucleo:} por el primer teorema de isomorfismo, $|\ker\phi|=|S_4|/|S_3|=24/6=4$.
Cada doble-trasposicion intercambia ambos elementos dentro de cada par de cada particion,
por lo que fija las tres particiones (aplicar la definicion de la accion a $(1\;2)(3\;4)$, $(1\;3)(2\;4)$ y $(1\;4)(2\;3)$).
Asi, $\ker\phi\supseteq V_4:=\{e,\;(1\;2)(3\;4),\;(1\;3)(2\;4),\;(1\;4)(2\;3)\}$;
como $|V_4|=4=|\ker\phi|$, concluimos $\ker\phi=V_4$.

\medskip
\textbf{Caso $n\ge 4$.}
Supongamos que existe $\phi\colon S_{n+1}\to S_n$ sobreyectivo.
Entonces $K:=\ker\phi$ es un subgrupo normal de $S_{n+1}$ de orden
\[
  |K|=\frac{|S_{n+1}|}{|S_n|}=\frac{(n+1)!}{n!}=n+1.
\]
Para $n\ge 4$ se tiene $n+1\ge 5$. Afirmamos que los unicos subgrupos normales de $S_{n+1}$ son
 $\{e\}$, $A_{n+1}$ y $S_{n+1}$: en efecto, si $N\trianglelefteq S_{n+1}$, entonces $N\cap A_{n+1}\trianglelefteq A_{n+1}$,
 y por el Lema 1 se tiene $N\cap A_{n+1}=\{e\}$ o $N\cap A_{n+1}=A_{n+1}$. Si $N\cap A_{n+1}=A_{n+1}$,
 entonces o bien $N=A_{n+1}$ o bien $N$ contiene una permutacion impar, y en este ultimo caso $N=S_{n+1}$.
 Si $N\cap A_{n+1}=\{e\}$ y $N$ contuviera una permutacion impar $\sigma$, entonces $\sigma^2\in A_{n+1}\cap N=\{e\}$,
 asi que $\sigma$ seria una trasposicion; como $N$ es normal, contendria todas las trasposiciones, y estas generan $S_{n+1}$,
 contradiccion. Por tanto, en este caso $N=\{e\}$. Asi, los unicos subgrupos normales de $S_{n+1}$ son:
\[
  \{e\}\;(\text{orden }1),\quad A_{n+1}\;\bigl(\text{orden }\tfrac{(n+1)!}{2}\bigr),\quad S_{n+1}\;\bigl(\text{orden }(n+1)!\bigr).
\]
Como $n\ge 4$, se tiene $n!\ge 24$. Si $(n+1)!=n+1$, entonces $n!=1$, imposible. Si $\tfrac{(n+1)!}{2}=n+1$, entonces $\tfrac{n!}{2}=1$, es decir, $n!=2$, imposible. Por tanto ninguno de estos ordenes coincide con $n+1$.
Por tanto no existe subgrupo normal de orden $n+1$ en $S_{n+1}$, y no puede existir tal homomorfismo sobreyectivo. \hfill$\square$

\problem{2}{%
Sea $G$ un grupo conmutativo y sea $C$ el conjunto de clases conjugadas de $G$, es decir, el conjunto de clases de equivalencia bajo la relacion
\[
  g_1 \sim g_2 \quad \text{si existe } h \in G \text{ tal que } g_2 = h g_1 h^{-1}.
\]
Dada una clase conjugada $c \in C$, defina
\[
  s_c := \sum_{g \in c} g
\]
como un elemento en $\CC[G]$. Demuestre que $s_c$ es un elemento central, es decir, que $s_c$ conmuta con todos los demas elementos del algebra de grupo.
}

\textbf{Respuesta.}
Basta demostrar que $h\,s_c = s_c\,h$ para todo $h\in G$,
pues los elementos $h\in G$ forman una base de $\CC[G]$ como $\CC$-espacio vectorial
y la igualdad se extiende por linealidad.

Equivalentemente, mostramos que $h\,s_c\,h^{-1}=s_c$.
Calculamos:
\[
  h\,s_c\,h^{-1}
  = h\Bigl(\sum_{g\in c}g\Bigr)h^{-1}
  = \sum_{g\in c}hgh^{-1}.
\]
Ahora bien, la conjugacion por $h$ es una biyeccion $G\to G$
que preserva cada clase conjugada: si $g\in c$, entonces $hgh^{-1}\sim g$
(pues $hgh^{-1}$ es conjugado de $g$), asi que $hgh^{-1}\in c$.
Por lo tanto la aplicacion $g\mapsto hgh^{-1}$ es una biyeccion $c\to c$
(es inyectiva por ser restriccion de una biyeccion, y $c$ es finito).

Al reindexar la suma con $g'=hgh^{-1}$, obtenemos
\[
  \sum_{g\in c}hgh^{-1}=\sum_{g'\in c}g'=s_c.
\]
Concluimos que $h\,s_c\,h^{-1}=s_c$, es decir $h\,s_c=s_c\,h$,
para todo $h\in G$. Esto muestra que $s_c\in Z(\CC[G])$. \hfill$\square$

\textit{Observacion:} el enunciado dice ``grupo conmutativo'', pero la demostracion
anterior vale para cualquier grupo finito $G$. Si $G$ es abeliano,
cada clase conjugada es un conjunto unitario $\{g\}$, $s_c=g$,
y el resultado se reduce a que $\CC[G]$ es conmutativo.

\problem{3}{%
Utilizando la terminologia del ejercicio anterior, demuestre que
\[
  Z(\CC[G]) = \operatorname{span}_{\CC}\langle s_c : c \in C \rangle.
\]
Donde, si $A$ es un anillo, $Z(A)$ denota el centro de $A$, es decir, el conjunto de elementos del algebra que conmutan con todos los demas.
}

\textbf{Respuesta.}

\textit{Contencion $\supseteq$.}
$Z(\CC[G])$ es un $\CC$-subespacio de $\CC[G]$, y por el Problema~2 cada $s_c$
pertenece a $Z(\CC[G])$; luego
$\operatorname{span}_{\CC}\langle s_c : c\in C\rangle\subseteq Z(\CC[G])$.

\medskip
\textit{Contencion $\subseteq$.}
Sea $z=\sum_{g\in G}\alpha_g\,g\in Z(\CC[G])$.
Para todo $h\in G$ se tiene $hzh^{-1}=z$, es decir
\[
  \sum_{g\in G}\alpha_g\,hgh^{-1}
  =\sum_{g\in G}\alpha_g\,g.
\]
Sustituyendo $g'=hgh^{-1}$ en el lado izquierdo
(biyeccion $G\to G$, con $g=h^{-1}g'h$):
\[
  \sum_{g'\in G}\alpha_{h^{-1}g'h}\,g'
  =\sum_{g'\in G}\alpha_{g'}\,g'.
\]
Como los elementos de $G$ son base de $\CC[G]$,
la igualdad coeficiente a coeficiente da
\[
  \alpha_{h^{-1}g'h}=\alpha_{g'}\qquad
  \text{para todos }g'\in G,\;h\in G.
\]
Esto dice que la funcion $g\mapsto\alpha_g$ es constante
en cada clase conjugada.
Denotando por $\alpha_c$ el valor comun para $g\in c$,
agrupamos la suma por clases conjugadas (que particionan $G$):
\[
  z=\sum_{c\in C}\;\sum_{g\in c}\alpha_g\,g
   =\sum_{c\in C}\alpha_c\sum_{g\in c}g
   =\sum_{c\in C}\alpha_c\,s_c
   \;\in\;\operatorname{span}_{\CC}\langle s_c:c\in C\rangle.
\]

De ambas contenciones,
$Z(\CC[G])=\operatorname{span}_{\CC}\langle s_c:c\in C\rangle$.
\hfill$\square$

\textit{Observacion:} los $s_c$ son $\CC$-linealmente independientes,
pues cada $g\in G$ aparece en exactamente un $s_c$ (el de su clase conjugada);
por tanto, en una relacion $\sum_c\lambda_c s_c=0$, el coeficiente de cualquier
$g\in c$ es $\lambda_c$, luego $\lambda_c=0$ para toda clase $c$.
Asi, $\{s_c:c\in C\}$ es una $\CC$-base de $Z(\CC[G])$ y
$\dim_\CC Z(\CC[G])=|C|$.

\problem{4}{%
Escriba $\CC[C_4]$ como suma directa de irreps.
}

\textbf{Respuesta.}
Sea $C_4=\langle r\mid r^4=e\rangle=\{e,r,r^2,r^3\}$ el grupo ciclico de orden~$4$.

Consideremos la representacion regular
\[
  \rho\colon C_4\to GL(\CC[C_4])\cong GL_4(\CC),
\]
donde $g\in C_4$ actua por multiplicacion a izquierda sobre $\CC[C_4]$.
En la base ordenada
\[
  \mathcal B=(e,r,r^2,r^3),
\]
la matriz de $\rho(r)$ es
\[
  P=
  \begin{pmatrix}
    0&0&0&1\\
    1&0&0&0\\
    0&1&0&0\\
    0&0&1&0
  \end{pmatrix},
\]
porque
\[
  r\cdot e=r,\qquad r\cdot r=r^2,\qquad r\cdot r^2=r^3,\qquad r\cdot r^3=e.
\]

Como $r^4=e$, se tiene $P^4=I$. Por tanto el polinomio minimo de $P$
divide a $x^4-1$, y sobre $\CC$ vale
\[
  x^4-1=(x-1)(x-i)(x+1)(x+i),
\]
que tiene cuatro raices distintas. Luego $P$ es diagonalizable y sus autovalores
pertenecen a $\{1,i,-1,-i\}$.

Definimos los vectores
\[
  v_0=\begin{pmatrix}1\\1\\1\\1\end{pmatrix},\qquad
  v_1=\begin{pmatrix}1\\-i\\-1\\i\end{pmatrix},\qquad
  v_2=\begin{pmatrix}1\\-1\\1\\-1\end{pmatrix},\qquad
  v_3=\begin{pmatrix}1\\i\\-1\\-i\end{pmatrix}.
\]
Se verifica directamente que
\[
  Pv_0=v_0,\qquad Pv_1=i\,v_1,\qquad Pv_2=-v_2,\qquad Pv_3=-i\,v_3.
\]
Asi, los subespacios $\CC v_0$, $\CC v_1$, $\CC v_2$ y $\CC v_3$
son invariantes por $C_4$ y de dimension~$1$; en particular, son irreducibles.

Si llamamos $\rho_k$ a la representacion $1$-dimensional dada por
\[
  \rho_k(r)=i^k,\qquad k=0,1,2,3,
\]
entonces
\[
  (\CC v_k,\rho|_{\CC v_k})\cong (\CC,\rho_k).
\]

\medskip
\textbf{Conclusion.}
Como representacion de $C_4$, la descomposicion de la representacion regular es
\[
  (\CC[C_4],\rho)
  \;\cong\; (\CC,\rho_0)\oplus(\CC,\rho_1)\oplus(\CC,\rho_2)\oplus(\CC,\rho_3).
\]
Equivalentemente,
\[
  \CC[C_4]
  \;=\; \CC v_0\oplus \CC v_1\oplus \CC v_2\oplus \CC v_3,
\]
y en cada sumando $\CC v_k$ la accion de $r$ es multiplicacion por $i^k$.
Esto es consistente con la dimension: $\dim\CC[C_4]=4=1+1+1+1$.
\hfill$\square$

\problem{5}{%
Escriba $\CC[S_3]$ como suma directa de irreps.
}

\begingroup
\small
\setlength{\parskip}{0.35em}
\setlength{\abovedisplayskip}{0.35em}
\setlength{\belowdisplayskip}{0.35em}
\setlength{\abovedisplayshortskip}{0.2em}
\setlength{\belowdisplayshortskip}{0.2em}
\setlength{\jot}{0.15em}

\textbf{Respuesta.}
Sea $P_3$ la representacion de permutacion de $S_3$ sobre $\CC^3$,
donde cada $\sigma\in S_3$ actua permutando la base canonica.
Consideremos los subespacios
\[
  L:=\CC(1,1,1),
  \qquad
  V:=\{(z_1,z_2,z_3)\in\CC^3:z_1+z_2+z_3=0\}.
\]
Ambos son $S_3$-invariantes: $L$ porque permutar coordenadas deja fijo al vector
$(1,1,1)$, y $V$ porque permutar coordenadas no cambia la suma $z_1+z_2+z_3$.
Ademas,
\[
  \CC^3=L\oplus V,
\]
porque $L\cap V=\{0\}$ y $\dim L+\dim V=1+2=3$.
Por tanto,
\[
  P_3\cong \mathbf 1\oplus V,
\]
donde $\mathbf 1$ es la representacion trivial y $V$ tiene dimension $2$.

\medskip
\textbf{Hecho (Corolario~4.3.15 de Steinberg~\cite[Cor.~4.3.15]{Steinberg}).}
Si $\chi$ es el caracter de una representacion compleja de un grupo finito $G$,
entonces la representacion es irreducible si y solo si
\[
  \langle \chi,\chi\rangle
  :=\frac{1}{|G|}\sum_{g\in G}\chi(g)\overline{\chi(g)}=1.
\]
Como los caracteres son constantes en clases conjugadas, la suma puede agruparse
por clases conjugadas.

Veamos que $V$ es irreducible. El caracter de $P_3$ se obtiene contando
cuantos vectores de la base canonica quedan fijos:
\[
  \chi_{P_3}(e)=3,
  \qquad
  \chi_{P_3}(\tau)=1\ \text{si }\tau\text{ es una trasposicion},
  \qquad
  \chi_{P_3}(c)=0\ \text{si }c\text{ es un }3\text{-ciclo}.
\]
Como $P_3\cong \mathbf 1\oplus V$, se sigue que
\[
  \chi_V(e)=2,
  \qquad
  \chi_V(\tau)=0,
  \qquad
  \chi_V(c)=-1.
\]
Entonces
\[
  \langle \chi_V,\chi_V\rangle
  =\frac{1}{|S_3|}\Bigl(1\cdot 2^2+3\cdot 0^2+2\cdot (-1)^2\Bigr)
  =\frac{1}{6}(4+0+2)=1.
\]
Por el criterio de irreducibilidad por caracteres, $V$ es irreducible.

Ahora determinamos todas las irreps de $S_3$.
Las clases conjugadas de $S_3$ son
\[
  \{e\},\qquad \{(1\;2),(1\;3),(2\;3)\},\qquad \{(1\;2\;3),(1\;3\;2)\}.
\]
Por tanto, sobre $\CC$ hay exactamente tres irreps, salvo equivalencia
(por el Teorema~4.4.8 de Steinberg~\cite[Thm.~4.4.8]{Steinberg}).
Las representaciones de dimension $1$ son exactamente dos: la trivial
y la signo. En efecto, toda representacion $\rho\colon S_3\to \CC^{\times}$
toma valores en $\CC^{\times}$, y este grupo es abeliano.
Por tanto, para cualesquiera $g,h\in S_3$ se tiene
$\rho(ghg^{-1}h^{-1})=1$, de modo que el subgrupo conmutador
$[S_3,S_3]$ queda contenido en $\ker\rho$.
Como $[S_3,S_3]=A_3$, se tiene $A_3\subseteq \ker\rho$. En particular,
si dos elementos de $S_3$ estan en la misma coclase modulo $A_3$,
entonces $\rho$ toma en ellos el mismo valor. Es decir, la aplicacion
$\rho$ solo depende de la coclase de $A_3$, y por tanto induce un
homomorfismo
\[
  \overline\rho\colon S_3/A_3 \longrightarrow \CC^{\times}.
\]
Como $S_3/A_3\cong C_2$, hay solo dos coclases, y toda representacion de
dimension $1$ de $S_3$ proviene de un homomorfismo $C_2\to \CC^{\times}$.
Estas corresponden precisamente a $\mathbf 1$ y $\operatorname{sgn}$.

Si las dimensiones de las irreps son $d_1,d_2,d_3$, entonces
(por el Corolario~4.4.5 de Steinberg~\cite[Cor.~4.4.5]{Steinberg})
\[
  d_1^2+d_2^2+d_3^2=|S_3|=6.
\]
Como ya tenemos dos irreps de dimension $1$, la tercera debe satisfacer
\[
  1^2+1^2+d_3^2=6,
\]
de donde $d_3=2$.
Esa tercera irrep es precisamente $V$.

Finalmente usamos el hecho general sobre la representacion regular
(Teorema~4.4.4 de Steinberg~\cite[Thm.~4.4.4]{Steinberg}):
si $G$ es finito y $W_1,\dots,W_r$ son sus irreps, de dimensiones
$d_1,\dots,d_r$, entonces
\[
  \CC[G]\cong d_1W_1\oplus d_2W_2\oplus \cdots \oplus d_rW_r.
\]
Aplicado a $G=S_3$, con irreps $\mathbf 1$, $\operatorname{sgn}$ y $V$,
obtenemos
\[
  \CC[S_3]\cong \mathbf 1\oplus \operatorname{sgn}\oplus V\oplus V.
\]

Esta es la descomposicion pedida. La comprobacion de dimensiones es
\[
  1+1+2+2=6=\dim_\CC \CC[S_3].
\]
\hfill$\square$
\endgroup

\problem{6}{%
Recuerde que $P_n$ es la representacion $\CC^n$ de $S_n$ donde cada elemento actua permutando los elementos de la base canonica. Verifique que $P_4$ se descompone como la suma directa de la irrep trivial y una irrep de dimension $3$.
}

\textbf{Respuesta.}
Consideremos los subespacios
\[
  L:=\CC(1,1,1,1),
  \qquad
  V:=\{(z_1,z_2,z_3,z_4)\in\CC^4:z_1+z_2+z_3+z_4=0\}.
\]
Ambos son $S_4$-invariantes: $L$ porque permutar coordenadas deja fijo
al vector $(1,1,1,1)$, y $V$ porque permutar coordenadas no cambia la suma
$z_1+z_2+z_3+z_4$.
Ademas, $L\cap V=\{0\}$ y $\dim L+\dim V=1+3=4=\dim\CC^4$, asi que
\[
  P_4\cong \mathbf 1\oplus V,
\]
donde $\mathbf 1$ denota la representacion trivial (sobre $L$).

\medskip
\textbf{Irreducibilidad de $V$.}
Usamos el criterio de irreducibilidad por caracteres
(Corolario~4.3.15 de Steinberg~\cite[Cor.~4.3.15]{Steinberg}):
una representacion con caracter $\chi$ es irreducible si y solo si
$\langle\chi,\chi\rangle=1$.

Las clases conjugadas de $S_4$ son cinco:
\[
\begin{array}{lcc}
  \text{Tipo de ciclo} & \text{Representante} & \text{Tamano} \\
  \hline
  (1,1,1,1) & e & 1 \\
  (2,1,1)   & (1\;2) & 6 \\
  (3,1)     & (1\;2\;3) & 8 \\
  (4)       & (1\;2\;3\;4) & 6 \\
  (2,2)     & (1\;2)(3\;4) & 3
\end{array}
\]
(Comprobacion: $1+6+8+6+3=24=|S_4|$.)

El caracter de la representacion de permutacion $P_4$ cuenta los puntos fijos
de cada permutacion:
\[
  \chi_{P_4}(e)=4,\quad
  \chi_{P_4}((1\;2))=2,\quad
  \chi_{P_4}((1\;2\;3))=1,\quad
  \chi_{P_4}((1\;2\;3\;4))=0,\quad
  \chi_{P_4}((1\;2)(3\;4))=0.
\]
Como $P_4\cong\mathbf 1\oplus V$ y $\chi_{\mathbf 1}\equiv 1$,
se tiene $\chi_V=\chi_{P_4}-\chi_{\mathbf 1}$:
\[
  \chi_V(e)=3,\quad
  \chi_V((1\;2))=1,\quad
  \chi_V((1\;2\;3))=0,\quad
  \chi_V((1\;2\;3\;4))=-1,\quad
  \chi_V((1\;2)(3\;4))=-1.
\]
Calculamos:
\[
  \langle\chi_V,\chi_V\rangle
  =\frac{1}{|S_4|}\sum_{g\in S_4}|\chi_V(g)|^2
  =\frac{1}{24}\bigl(1\cdot 3^2+6\cdot 1^2+8\cdot 0^2+6\cdot(-1)^2+3\cdot(-1)^2\bigr)
  =\frac{9+6+0+6+3}{24}=\frac{24}{24}=1.
\]
Por el criterio de irreducibilidad, $V$ es una representacion irreducible de dimension~$3$.

\medskip
\textbf{Conclusion.}
\[
  P_4\;\cong\;\mathbf 1\;\oplus\; V,
\]
donde $\mathbf 1$ es la irrep trivial (dimension $1$) y $V$ es una irrep de dimension $3$.
\hfill$\square$

\medskip
\textit{Observacion (caso general).}
La descomposicion $P_n\cong\mathbf 1\oplus V_n$, con
$V_n:=\{(z_1,\dots,z_n)\in\CC^n:z_1+\cdots+z_n=0\}$
irreducible de dimension $n-1$, vale para todo $n\ge 2$.
La descomposicion como suma directa es identica al caso $n=4$:
$L=\CC(1,\dots,1)$ y $V_n=\ker(\text{suma})$ son $S_n$-invariantes,
$L\cap V_n=\{0\}$, y $\dim L+\dim V_n=n$.
Para la irreducibilidad, el caracter de $P_n$ es
$\chi_{P_n}(\sigma)=|\operatorname{Fix}(\sigma)|$,
asi que $\chi_{V_n}(\sigma)=|\operatorname{Fix}(\sigma)|-1$.
El calculo de $\langle\chi_{V_n},\chi_{V_n}\rangle$ no requiere
enumerar clases conjugadas: basta usar
\[
  \sum_{\sigma\in S_n}|\operatorname{Fix}(\sigma)|=n!,
  \qquad
  \sum_{\sigma\in S_n}|\operatorname{Fix}(\sigma)|^2=2\cdot n!,
\]
que se demuestran por conteo directo
(la primera suma cuenta pares $(i,\sigma)$ con $\sigma(i)=i$,
contribucion $n\cdot(n-1)!=n!$;
la segunda cuenta pares $((i,j),\sigma)$ con $\sigma(i)=i$ y $\sigma(j)=j$,
contribucion $n\cdot(n-1)!+n(n-1)\cdot(n-2)!=2\cdot n!$).
Entonces
\[
  \langle\chi_{V_n},\chi_{V_n}\rangle
  =\frac{1}{n!}\sum_{\sigma\in S_n}(|\operatorname{Fix}(\sigma)|-1)^2
  =\frac{2n!-2n!+n!}{n!}=1.
\]
La representacion $V_n$ se llama la \emph{representacion estandar} de $S_n$.

\begin{thebibliography}{9}

\bibitem{Steinberg}
Benjamin Steinberg,
\textit{Representation Theory of Finite Groups: An Introductory Approach},
Universitext, Springer, 2012.

\end{thebibliography}

\end{document}
